\subsection{E4: closedness of the numerical locus without projectivity}

\paragraph{Theorem E4.T.}
Let $f:X\to S$ be a proper flat morphism of finite presentation with
geometrically integral fibres. The base scheme $S$ is arbitrary. For every
invertible sheaf $L$ on $X$, the set
\[
 S^\tau(L)=\{s\in S:L_{\bar s}\text{ is numerically trivial}\}
                                                        \tag{E4.T.1}
\]
is open and closed. Equivalently, the open subgroup algebraic space
$\operatorname{Pic}^{\tau}_{X/S}$ is closed in
$\operatorname{Pic}_{X/S}$. Thus projectivity is unnecessary in the
question ``projectif est-il n\'ecessaire ?'' on 57:314. Here numerical
triviality and tau-equivalence agree for proper schemes over a field;
see \cite[Remark~6.14, pp.~57--58]{Kleiman2005}.

\paragraph{1. The trait assertion.}
Let $R$ be a discrete valuation ring, with fraction field $K$, and let
$X/R$ satisfy the hypotheses. Suppose $L_K$ is numerically trivial.
We prove that $L_k$ is numerically trivial. Replacing $R$ by its completion
preserves the special fibre and the hypotheses; numerical triviality of
the generic bundle is preserved under field extension. We may therefore
assume $R$ complete, hence excellent. Flatness and integrality of the generic
fibre make $X$ integral. If $\xi$ is the generic point of its reduced
integral special fibre, $\mathcal O_{X,\xi}$ is a one-dimensional local
domain with maximal ideal generated by a uniformizer of $R$. It is a DVR.

Zheng's Galois alteration lemma
\cite[Lemma~3.8, pp.~313--314]{BCEZheng2009} gives a finite extension of
complete traits $R\subset R'$, a regular integral scheme $Y$ projective
over $R'$, and a proper generically finite map
\[
 \pi:Y\longrightarrow X
\]
equivariant for a finite group $G$ acting trivially on $X$. Moreover
$k(Y)^G/k(X)$ is purely inseparable. This is precisely the definition of
Galois alteration used on p.~313 of that source. To apply the lemma take
the input group trivial and the dense open $X_K$. Its compactification
$Y$ is the projective end of a plurinodal fibration over $R'$. The source
$X'$ of the open immersion in the general statement is proper over $X$,
hence over $R$. Since $R'/R$ is separated, $X'$ is proper over $R'$ as
well. Its map into the separated projective end is therefore proper.
The dense open immersion is consequently an isomorphism. The lemma also makes the
boundary a strict normal-crossings divisor. Only regularity of $Y$ and
the intersection theory below are needed: regularity of a geometric
special fibre is not assumed.

Put $M=\pi^*L$ and write
\[
 Y_{k'}=\sum_{i=1}^r m_iY_i,\qquad m_i>0.                \tag{E4.T.2}
\]
The special fibre is connected. Indeed $H^0(Y,\mathcal O_Y)$ is finite
over the complete local ring $R'$ and is a domain, hence a local ring;
Stein factorization then gives connectedness. Each $Y_i$ is a Cartier
divisor on the regular $Y$. Let $d=\dim Y_{K'}$. The case $d=0$ is
immediate, since every line bundle on a proper zero-dimensional scheme
is numerically trivial. Assume $d\geq1$, and choose a relatively ample
line bundle $H$ on $Y$.

\paragraph{2. The vertical correction.}
Define the rational symmetric matrix $A$ and vector $b$ by intersection
numbers on the special fibre:
\[
 A_{ij}=(Y_i\cdot Y_j\cdot H^{d-1}),\qquad
 b_i=(M\cdot Y_i\cdot H^{d-1}).                          \tag{E4.T.3}
\]
For $i\ne j$ one has $A_{ij}\geq0$, with strict inequality exactly when
$Y_i$ and $Y_j$ meet. This follows because their proper intersection is
an effective cycle and $H$ is ample on the projective special fibre.
The relation $\sum_jm_jY_j=\operatorname{div}(\varpi')$ gives $Am=0$.
The intersection graph is connected. Consequently
\[
 v^tAv=-\sum_{i<j}A_{ij}m_im_j
       \left(\frac{v_i}{m_i}-\frac{v_j}{m_j}\right)^2,
 \qquad \ker A=\mathbf Qm.                              \tag{E4.T.4}
\]
Constancy of intersection numbers in a flat projective family gives
\[
 \sum_i m_i b_i=(M_{K'}\cdot H_{K'}^{d-1})=0.
\]
Thus there is $a\in\mathbf Q^r$ with $Aa=-b$. Set
\[
 D=\sum_i a_iY_i,\qquad N=M+D\quad\text{in }\operatorname{Pic}(Y)_\mathbf Q.
                                                               \tag{E4.T.5}
\]
By construction $(N|_{Y_i}\cdot H|_{Y_i}^{d-1})=0$ for every $i$.
For $d=1$ this already says that $N|_{Y_i}$ is numerically trivial.
For $d\geq2$, the Hodge index theorem gives
\[
 (N|_{Y_i}^{\,2}\cdot H|_{Y_i}^{d-2})\leq0.             \tag{E4.T.6}
\]
Equality implies numerical triviality. The higher-dimensional version
follows by cutting with sufficiently ample integral hyperplanes and the
surface Hodge theorem; the equality implication is
\cite[Theorem~6.3(h), proof pp.~53--54, and Theorem~B.27, p.~79]{Kleiman2005}.
These statements apply to integral projective schemes, with no smoothness
hypothesis. Over a nonclosed field, pass to an algebraic closure: the
geometric components of an integral scheme are conjugate, and the
intersection numbers of descended bundles are equal on these components
up to their common positive multiplicity. The degree-zero condition,
the inequality, and its equality conclusion therefore give the same
statement over the original field. Denominators in $N$ are cleared before
applying the line-bundle theorem.

The same flat-family intersection calculation gives
\[
 \sum_i m_i(N|_{Y_i}^{\,2}\cdot H|_{Y_i}^{d-2})
       =(M_{K'}^{\,2}\cdot H_{K'}^{d-2})=0.              \tag{E4.T.7}
\]
Every summand is nonpositive and every $m_i$ is positive. Each is therefore
zero, and $N$ is numerically trivial on every special component, hence
on the whole special fibre. It is also numerically trivial on the generic
fibre. This construction gives a rational vertical correction; it does
not assert that an arbitrary generic line bundle extends integrally.

\paragraph{3. Galois symmetry removes the correction.}
Replace $D$ by $|G|^{-1}\sum_{g\in G}g^*D$. This is a rational average,
so no invertibility assumption on $|G|$ in $R$ is used. Since $M$ is pulled
back from $X$, it is $G$-invariant, and numerical triviality is preserved
by each (possibly semilinear over $R'$) automorphism. Thus $M+D$ remains
numerically trivial on both fibres.

Let $Z$ be the normalization of $X$ in $k(Y)$. It is finite over $X$ by
excellence. Normality of $Y$ factors $\pi$ as
\[
 Y\xrightarrow{h}Z\longrightarrow X,
\]
where $h$ is proper birational between normal schemes. The finite group
$G$ acts transitively on the prime divisors of $Z$ above $\xi$.
Here is the precise valuation argument. The DVR $\mathcal O_{X,\xi}$ has
a unique extension to the purely inseparable field $k(Y)^G$; the field
extension $k(Y)/k(Y)^G$ is Galois with group the faithful quotient of $G$,
and its Galois group permutes transitively the extensions of that valuation.
The same transitivity holds for their strict transforms on $Y$, since
a proper birational map to a normal scheme is an isomorphism at the
generic point of every target prime divisor.

These are exactly the components $Y_i$ that dominate $X_k$. Indeed a
vertical prime of the finite scheme $Z\to X$ has dimension $d$, hence
maps onto the integral $d$-dimensional special fibre. Their multiplicities
$m_i$ in (E4.T.2) are constant on the orbit: $G$ fixes the uniformizer of
$R$, whose divisor is the ramification index of $R'/R$ times (E4.T.2).
The averaged coefficients $a_i$ are constant on this orbit as well.
Consequently $a_i/m_i$ has a single value $c$ on every nonexceptional
vertical component. Subtract $c\operatorname{div}(\varpi')$ from $D$.
This is a principal rational divisor, so it preserves all numerical
conclusions and makes $D$ supported on $h$-exceptional prime divisors.

For any curve contracted by $h$, the degrees of $M$ and of $N=M+D$ are
zero. This includes curves over generic-fibre points, since there $D=0$
and $N=M$ is numerically trivial, as well as the special-fibre curves
treated above. Hence both $D$ and $-D$ are $h$-nef. The negativity lemma for proper
birational morphisms of normal schemes
\cite[Lemma~11.60, p.~436]{KollarFamilies} applied to both signs makes
$D\geq0$ and $D\leq0$, because $h_*D=0$. Thus $D=0$.
We conclude that $M|_{Y_{k'}}=\pi^*L|_{Y_{k'}}$ is numerically trivial.
Numerical triviality is detected by proper surjective pullback: for an
integral curve $C\subset X_k$, take a closed point of the generic fibre
of $Y_{k'}\times_{X_k}C\to C$ and close it to an integral curve $C'$.
The map $C'\to C$ is generically finite of positive degree, and the
projection formula gives
$0=\deg(M|_{C'})=[k(C'):k(C)]\deg(L|_C)$.
This proves the trait assertion.

\paragraph{4. From traits to an arbitrary base.}
First suppose $S$ is locally noetherian. The set (E4.T.1) is open by
the proper-flat Picard tau theorem
\cite[Theorem~2.1]{LaurentSchroer2024}; proper geometrically integral
fibres give universal $f_*\mathcal O_X=\mathcal O_S$.
Every specialization in a locally noetherian scheme is realized by a
map from the spectrum of a DVR (Stacks, Lemma~28.5.10, tag~054F).
The trait assertion shows that (E4.T.1) is stable under specialization.
Locally on the noetherian base it is constructible, being open; a
constructible specialization-stable subset is closed. Thus it is open
and closed.

For arbitrary $S$ the assertion is local, so take $S=\operatorname{Spec}A$.
Noetherian approximation descends $X/S$ and $L$ to some scheme
$S_i=\operatorname{Spec}A_i$ of finite type over $\mathbf Z$, with $X_i/S_i$
proper, flat and finitely presented. The needed exact limit results are
Stacks, Lemmas~32.10.1, 32.13.1, and 32.8.7 (tag~04AI), together with
finite-presentation descent of an invertible module. To ensure geometric
integrality of every fibre, let $E_i\subset S_i$ be the locus where the
fibre is geometrically reduced and geometrically irreducible. It is
constructible by Lemmas~37.26.5 and 37.27.7 (the latter tag~055B), and
$S\to S_i$ has image in $E_i$. Lemma~32.4.10, tag~05F4, shows that for
a sufficiently large index the entire $S_j\to S_i$ has image in $E_i$.
Replacing the model by $S_j$ therefore gives proper flat geometrically
integral fibres, while retaining $L_j$.
The cited limit and constructibility sources are collected in
\cite{BCEStacksTauDescent}.

The already proved noetherian assertion says $S_j^\tau(L_j)$ is open and
closed. Numerical triviality is invariant under field extension, so
$S^\tau(L)$ is exactly its inverse image. This proves (E4.T.1) for every
base. Finally, an fppf-local representative of a Picard section is an
actual invertible sheaf; applying (E4.T.1) on such a cover and descending
the open-and-closed locus proves the asserted closed immersion of the
tau subfunctor into the full Picard algebraic space. No schematicity
of the full Picard functor over $S$ is assumed or deduced.
